\documentclass[10pt]{article}
\usepackage[T1]{fontenc}
\usepackage[utf8]{inputenc}
\usepackage[french]{babel}
\usepackage{amsmath,amssymb,mathrsfs,mathpazo}
\usepackage{geometry}
\geometry{a4paper,margin=2cm}
\usepackage{graphicx,array,tabularx,tabularray,colortbl,multirow,multicol,booktabs,enumitem,pifont,ulem,fancybox,diagbox,cancel}
\usepackage{pstricks,pstricks-add,pst-plot,pst-tree,pst-node,pst-text,pst-func,pst-math,pst-bspline,pst-3dplot,pst-eucl,pst-solides3d}
\let\pscalculate\undefined
\usepackage{pst-bezier}
\usepackage{tikz,pgfplots,tkz-tab}
\usetikzlibrary{arrows,arrows.meta,patterns,calc,positioning,trees,decorations.markings,shapes.geometric,shapes.misc,shapes.arrows,intersections,decorations.pathmorphing,decorations.pathreplacing,backgrounds,shadows,babel}
\pgfplotsset{compat=1.18}
\usepgfplotslibrary{fillbetween}
\usepackage[np]{numprint}
\usepackage{listings,makecell,pas-tableur,twoopt}
\let\barre\undefined
\usepackage[pstricks]{tablvar}
\definecolor{darkgreen}{RGB}{0,100,0}
\providecommand{\red}{\color{red}}
\providecommand{\blue}{\color{blue}}
\providecommand{\green}{\color{green}}
\providecommand{\black}{\color{black}}
\providecommand{\white}{\color{white}}
\providecommand{\gray}{\color{gray}}
\providecommand{\cyan}{\color{cyan}}
\providecommand{\magenta}{\color{magenta}}
\providecommand{\yellow}{\color{yellow}}
\providecommand{\violet}{\color{violet}}
\providecommand{\orange}{\color{orange}}
\providecommand{\euro}{\texteuro}
\providecommand{\EUR}{\texteuro}
\providecommand{\ieme}{\textsuperscript{e}}
\providecommand{\degres}{\ensuremath{^\circ}}
\providecommand{\degree}{\ensuremath{^\circ}}
\providecommand{\vv}[1]{\overrightarrow{#1}}
\providecommand{\Coord}[1]{#1}
\providecommand{\decosix}{}
\providecommand{\starredbullet}{\ensuremath{\bullet}}
\pagestyle{empty}
\setlength{\parindent}{0pt}
\usepackage{adjustbox}\DeclareUnicodeCharacter{2236}{:}\DeclareUnicodeCharacter{2190}{\ensuremath{\leftarrow}}
\providecommand{\R}{\mathbb{R}}
\providecommand{\N}{\mathbb{N}}
\providecommand{\Z}{\mathbb{Z}}
\providecommand{\Q}{\mathbb{Q}}
\providecommand{\D}{\mathbb{D}}
\providecommand{\C}{\mathbb{C}}
\providecommand{\e}{{\mathrm{e}}}
\providecommand{\vect}[1]{\overrightarrow{#1}}
\providecommand{\vectt}[1]{\overrightarrow{\mathrm{#1}}}
\providecommand{\vvt}[1]{\overrightarrow{\mathrm{#1}}}
\providecommand{\np}[2][]{{#2~#1}}
\providecommand{\vv}[1]{{\overrightarrow{#1}}}
\renewcommand{\d}{{\mathrm{d}}}
\providecommand{\degres}{{^\circ}}
\providecommand{\degre}{{^\circ}}
\providecommand{\blue}{\color{blue}}
\providecommand{\red}{\color{red}}
\providecommand{\enskip}{\quad}
\providecommand{\llbracket}{\lbrack\!\lbrack}
\providecommand{\rrbracket}{\rbrack\!\rbrack}
\providecommand{\euro}{{\text{€}}}
\providecommand{\no}{{\text{n°}}}
\providecommand{\interval}[2]{{[#1;#2]}}
\providecommand{\intervFO}[2]{{[#1;#2[}}
\providecommand{\dotfill}{\ldots}

\providecommand{\ds}{\displaystyle}
\providecommand{\textding}[1]{#1}
\providecommand{\Oij}{\(\mathrm{O};\vec{\imath},\vec{\jmath}\)}
\providecommand{\Oijk}{\(\mathrm{O};\vec{\imath},\vec{\jmath},\vec{k}\)}
\providecommand{\Ouv}{\(\mathrm{O};\vec{u},\vec{v}\)}
\providecommand{\intervFF}[2]{\ensuremath{[#1;#2]}}
\renewcommand{\euro}{\ifmmode\text{\texteuro}\else\texteuro\fi}
\DeclareUnicodeCharacter{00B0}{\ifmmode^{\circ}\else\textdegree\fi}
\DeclareUnicodeCharacter{20AC}{\ifmmode\text{\texteuro}\else\texteuro\fi}
% Énoncés : Éducation nationale. Transcription et corrigés : APMEP, Denis Vergès et contributeurs. Découpage et classement : CoopMaths, Sébastien Lozano, Éric Elter et contributeurs. Adaptation de la présentation : Maths Explorer.
% CC BY-SA : https://coopmaths.fr/www/about/
% https://coopmaths.fr/alea/static/bac/2026/tex/bac_2026_06_sujet2_asie_4_cor.tex
\begin{document}
\noindent\textit{APMEP / CoopMaths. CC BY-SA. Présentation adaptée par Maths Explorer.}\par\medskip
\begin{enumerate}
\item \textbf{Affirmation 1 : Fausse.}

Calculons la dérivée seconde de $f(x) = (-0,5x + 3)^5$. (on rappelle que $(u^n)'=nu^{n-1}u'$)
\[
\begin{array}{rcl}
f'(x) &=& 5 \times (-0,5) \times (-0,5x + 3)^4 = -2.5(-0,5x + 3)^4 \\
f''(x) &=& -2,5 \times 4 \times (-0,5) \times (-0,5x + 3)^3 = 5(-0,5x + 3)^3
\end{array}
\]
Pour $x = 10$, on a $f''(10) = 5(-5 + 3)^3 = 5(-8) = -40 < 0$. La dérivée seconde n'est pas positive sur tout $\mathbb{R}$, la fonction n'est donc pas convexe sur $\mathbb{R}$.

\item \textbf{Affirmation 2 : Vraie.} 

Il y a 4 multiples de 8 parmi les 32 jetons $\{8, 16, 24, 32\}$.
Le nombre de tirages contenant au moins un multiple de 8 est le nombre total de tirages moins le nombre de tirages ne contenant aucun multiple de 8 (tirage parmi les 28 autres jetons).
\[
\begin{array}{rcl}
\text{Nombre total} &=& \binom{32}{5} = 201376 \\
\text{Sans multiple de 8} &=& \binom{28}{5} = \np{98280} \\
\text{Au moins un} &=& \np{201376} - \np{98280} = \np{103096}
\end{array}
\]

\item \textbf{Affirmation 3 : Fausse.} Complétons l'arbre fourni par l'énoncé.

\begin{center}
\begin{tikzpicture}[
grow=right,
level distance=4cm,
level 1/.style={sibling distance=4cm},
level 2/.style={sibling distance=1.5cm},
every node/.style={ circle, inner sep=2pt}
]
\node[coordinate] {}
child { node {$\overline{A}$}  
 child { node {$\overline{B}$} edge from parent node[below] {$\frac7{10}$}}
child { node {B} edge from parent node[above] {$\frac3{10}$ }}
edge from parent node[below] {$\frac35$}
}
child { node {$A$}
child { node {$\overline{B}$} edge from parent node[below] {$\frac34$}}
child { node {B} edge from parent node[above] {$\frac14$ }}
edge from parent node[above] {$\frac25$}};
\end{tikzpicture}
\end{center}
On peut dans un premier temps calculer $p(B)$ par la formule de probabilités totales:\[
p(B)= p(A \cap B) + p\left(\overline{A} \cap B\right) = \frac25 \times \frac14 + \frac35 \times \frac3{10} =\frac1{10}+\frac9{50}= \frac{14}{50}
\]
On en déduit que: \[
p_B\left(\overline{A}\right)=\frac{p\left(\overline{A} \cap B\right)}{p(B)} =\frac{\frac9{50}}{\frac{14}{50}} = \frac9{14}
\]


\item \textbf{Affirmation 4 : Vraie.} \\
Vérifions si $h(x) = \mathrm{e}^{-x}\sin(x)$ vérifie $h'(x) + h(x) = \mathrm{e}^{-x} \cos(x)$.
\[
\begin{array}{rcl}
h'(x) &=& -\mathrm{e}^{-x}\sin(x) + \mathrm{e}^{-x}\cos(x) \\
h'(x) + h(x) &=& \left( -\mathrm{e}^{-x}\sin(x) + \mathrm{e}^{-x}\cos(x) \right) + \mathrm{e}^{-x}\sin(x) \\
h'(x) + h(x) &=& \mathrm{e}^{-x}\cos(x)
\end{array}
\]
La fonction $h$ est bien une solution particulière.

\item \textbf{Affirmation 5 : Fausse.} 

Les solutions de $(E)$ sont la somme de la solution particulière $h(x)$ et des solutions de l'équation homogène associée $y' + y = 0$.
Les solutions de l'équation homogène sont de la forme $C\mathrm{e}^{-x}$. 
L'ensemble des solutions de $(E)$ est donc constitué des fonctions de la forme $y(x) = C\mathrm{e}^{-x} + \mathrm{e}^{-x}\sin(x)$, ce qui est différent de l'expression proposée.
\end{enumerate}
\label{LastPage}

\end{document}
